Напомним следующие понятия, чтобы было проще ориентироваться

\begin{definition} Классом \bpp (bounded-error probabilistic polynomial) называется класс языков $A$, для которых существует полиномиальный в худшем случае вероятностный алгоритм $V$, такой что:
\begin{itemize}
    \item если $x \in A$, то $\Pr_r[V(x, r) = 1] \geq \frac{2}{3}$;
    \item если $x \not\in A$, то $\Pr_r[V(x, r) = 1] \leq \frac{1}{3}$.
\end{itemize}
\end{definition}


\begin{definition} Пусть $\alpha (n), \beta(n)$ -- две функции из $\N$ в $[0,1]$, такие что $0 \leq \alpha(n) < \beta(n) \leq 1$ при всех $n$. Классом $\bpp_{\alpha(n), \beta(n)}$ будем называть $\bpp$, для которого пороговые вероятности таковы:
\begin{itemize}
    \item если $x \in A$, то $\Pr_r[V(x, r) = 1] \geq \beta(|x|)$;
    \item если $x \not\in A$, то $\Pr_r[V(x, r) = 1] \leq \alpha(|x|)$.
\end{itemize}
\end{definition}

Таким образом, $\alpha(|x|)$ является верхней оценкой на вероятность $\alpha$-ошибки, а $1-\beta(|x|)$ -- верхней оценкой на вероятность $\beta$-ошибки.
В новых обозначения можно записать классы: $\bpp = \bpp_{\frac{1}{3}, \frac{2}{3}}, \rp = \bpp_{0, \frac{1}{2}}, \corp = \bpp_{\frac{1}{2}, 1}$.

\begin{theorem}
Пусть $c, d \in \Z$. Тогда $\bpp_{0, \frac{1}{n^c}} = \bpp_{0, 1 - \frac{1}{2^{n^d}}} = \rp$ и $\bpp_{1 - \frac{1}{n^c}, 1} = \bpp_{\frac{1}{2^{n^d}}, 1} = \corp$
\end{theorem}

\begin{proof}
Оба утверждения доказываются абсолютно аналогично, поэтому для примера докажем только первое. Новый алгоритм будет таким: запускаем старый алгоритм $k$ раз с независимыми случайными битами. Если хотя бы один раз получилась единица, возвратим единицу, иначе ноль. (То есть возвратим дизъюнкцию результатов). Если на самом деле ответ <<да>>, то вероятность того, что все $k$ раз алгоритм вернет $0$, равна $(1 - \frac{1}{n^c})^k$. В силу второго замечательного предела это примерно равно $e^{-\frac{k}{n^c}}$. Таким образом, если взять $k = n^{c+d}$, то вероятность ошибки не превысит $\frac{1}{2^{n^d}}$, ЧТД. Равенство $\rp$ доказывается вложениями $\bpp_{0, 1 - \frac{1}{2^{n^d}}} \subset \rp \subset \bpp_{0, \frac{1}{n^c}}$
\end{proof}

\begin{theorem} Пусть $c, d \in \Z$, $\tau \in (0, 1)$. Тогда $\bpp_{\tau - \frac{1}{n^c}, \tau + \frac{1}{n^c}} = \bpp_{\frac{1}{2^{n^d}}, 1 - \frac{1}{2^{n^d}}}$. В частности $\forall \eps \in \brackets{0, \frac{1}{2}} \hookrightarrow \bpp_{\eps, 1 - \eps} = \bpp$.
\end{theorem}

\begin{proof} Запустим алгоритм много раз (то есть проведём амплификацию), а затем сравним долю успехов с $\tau$. Если она больше $\tau$, то вернём 1, иначе — 0. Заметим, что нам необходимо делать полиномиальную амплификацию (то есть запускать алгоритм не более, чем полином раз).

Для примера, пусть вероятность успеха при одном запуске $\geq \tau + \frac{1}{n^c}$. Обозначим за $\xi_i$ случайную величину~--- ответ, который машина возвращает на $i$-м запуске. Тогда, очевидно, $\xi_i \sim \text{Bern}\brackets{\tau + \frac{1}{n^c}}$. Тогда нам надо подобрать такое $k$, что $\sum\limits_{i = 1}^k \xi_i \geq \tau \cdot k$.

А именно, при полиномиальном $k$ эта вероятность должна быть экспоненциально близка к единице. Интуитивно это объясняется так. Рассмотрим $S_k = \frac{\sum\limits_{i = 1}^k \xi_i}{k}$ и применим к ней ЦПТ:

$$
S_K = \frac{\sum\limits_{i = 1}^k \xi_i}{k} \sim \mathcal{N}\brackets{\tau + \frac{1}{n^c}, O\brackets{\frac{1}{k}}}
$$

Тогда через $k = n^{2c} \cdot L$ запусков получим, что $\tau$ отстоит от среднего примерно на $\sqrt{L}$ стандартных отклонений, поэтому вероятность того, что $S_k < \tau$, экспоненциальная ($\exp(-L)$, из свойств нормального распределения), что и требовалось.
\end{proof}